\documentclass[12pt,a4paper,oneside]{report}
% nagyon sok kép esetén meggyorsítható a fordítás a draft móddal
% \documentclass[12pt,a4paper,oneside,draft]{report}
% ekkor a képek nem renderelődnek ki, csak placeholder lesz mérethelyesen
\usepackage[utf8]{inputenc}
\usepackage[magyar]{babel}
\usepackage[T1]{fontenc}
%\usepackage{amsmath}
%\usepackage{amsfonts}
%\usepackage{amssymb}
\usepackage{graphics} % grafikus elemek, képek berakásához
\usepackage{epsfig} % eps importáláshoz
%\usepackage{listings}
%\usepackage{sectsty}
\usepackage{enumerate}
\usepackage{lastpage}
\usepackage{setspace}
\usepackage{hyperref} % PDF hivatkozásokhoz kell
\usepackage[hang]{caption}
\usepackage{titling} % a title, author parancsok szabad használatához

% a TikZ rajzoló modul, és a kapcsolási rajz készítő modul, ha kell
\usepackage{tikz}
\usetikzlibrary{automata,positioning}
\usetikzlibrary{shapes,arrows,fit,calc,positioning,automata}
%\usepackage{circuitikz}

% az A4 oldal margóinak és méreteinek beállítása
\usepackage[left=25mm,right=25mm,top=20mm,bottom=25mm]{geometry}\pagestyle{plain}

% A sorköz távolság beállítása
% egyszeres sorköz
\singlespacing
% 1,5 sorköz
% \onehalfspacing

% A hivatkozások, és linkek átállítása alapértelmezett színre, fekete-fehér nyomtatáshoz optimalizálva
\hypersetup
{
  	colorlinks,
  	citecolor=black,
 	linkcolor=black,
  	urlcolor=black
}

\title{Mátrix4 dokumentáció} 

\author{Simonyi Károly Szakkollégium - Schönherz Elektronikai Műhely}

\begin{document}
\begin{titlepage}

\begin{figure}
\centering
\includegraphics[width=100mm,keepaspectratio]{images/logos/sem.png}
\end{figure}

\centering

\vspace{5mm}
\includegraphics[width=40mm,keepaspectratio]{images/logos/simonyi.png}  \\
\textbf{Villamosmérnöki és Informatikai Kar}\\
\textbf{Budapesti Műszaki és Gazdaságtudományi Egyetem}\\
\vspace{50mm}
\Huge
{\let\newpage\relax\maketitle}

\end{titlepage}
\chapter*{Előszó}

A 2017-es mátrix felszereléskor már nem volt tartalék MUEB. A jövőre való tekintettel újraterveztük a Mátrix hardveres részét, hogy egyrsézt továbbra is életben tudjuk tartani a hagyománynak számító eseményt, másrészt a jövő generációinak megkönnyítsük a karbantartási és felszerelési feladatokat.

\par
Résztvevők:
\begin{enumerate} %TODO pontosítás, bővítés
 \item Szabó Benjámin (SEM)  NYÁK tervezés, és dobozolás
 \item Kiss Ádám (SEM)
 \item Sweidán Omár (SEM) MUEB teszt szoftver
\end{enumerate}

\tableofcontents
\chapter{Koncepció}
%TODO szószedet írni, pl MUEB szó feloladása

A kitűzött cél a MUEB-ek minél nagyobb hibatűrés, és debugolhatósága. Ezen kívül a konfigurálást is minél inkább automatizálni kell.
\par
Ennek érdekében a következő célokat tűztük ki:
\begin{enumerate}
 \item Visszajelző LED-ek a MUEB-en. Az egyértelműség kedvéért legyen rászitázva a LED-ek mellé, hogy melyik mit csinál.
 \item A MUEB automata azonosítsa be, hogy melyik szobában van, ne kelljen jumperelni.
 \item Az áramkör legyen zavartűrő, és ne bocsásson ki magából zajt
 \item Ne lehessen véletlen zárlatot csinálni a NYÁK-on
 \item Ne lehessen elcsúsztatva bedugni a csatlakozókat
 \item Egyértelmű legyen, hogy hova kell dugni a tápcsatlakozót
 \item Ne legyen kritikus a kábelek csatlakoztatásának a sorrendje.
 \item Ergonómikus dobozban legyen
 \item Az elkészült kódokat tároljuk könnyen azonosítható helyen, valamint készüljön hozzájuk jól érthető dokumentáció a jövőbeli változtatások megkönnyítése érdekében.
 \item Szerver beüzemelése nélkül is lehessen animációt lejátszani a MUEB-en
 \item A MUEB működése stabil legyen, az esetleges hibák pedig következetesek.
 \item A MUEB illeszkedjen az SCHNET-hez, azon semmilyen módosítást ne kelljen végrehajtani az esemény idejére.
\end{enumerate}
Ami nem célunk:
\begin{enumerate}
 \item Más épületekbe hordozható hardver készítése
\end{enumerate}

\chapter{Hardver}

\section{Visszajelző szervek}
A NYÁK-on 8 db LED található.
A szoftverből vezérelt LED-ek funkciójáról részletesebben a %TODO
fejezetben olvashatnak.

\subsection{12V feliratú LED}
A 12V feliratú LED jelzi, hogy a csatlakoztatott külső tápegység feszültséget ad ki magából, illetve ép az áramkör bementén lévő 1,5A-es biztosíték.
\subsection{3.3V fleiratú LED}
A 3.3V feliratú LED jelzi, hogy a %TODO
%TODO ablak ledek

\subsection{Heartbeat LED}
A mikrokontroller épségét jelzi vissza. A szoftver villogtatja másodperc nagyságú periódusidővel.
\subsection{DHCP LED}
A DHCP-vel kapcsolatos 



\chapter{A kijelzők vezérlőpanelje}

\section{Koncepció}
Egy kijelző 4 darab pixelt tartalmaz. A panel egy szalagkábelen keresztül csatlakozik a MUEB-hez, azzal UART-on ommunikálk.
A két ablkot egy jumperelés konfigurálja ``ellentétes'' funkcióra.

\subsection{Pixelek}
Minden pixel 9 darab szubpixelből áll össze. A szubpixeleken belül három azanos színű LED van sorba kötve, majd ezek a blokkok párhuzamosan.
%todo a táblán lévő ledek elhelyezkedése
%todo szalagkábelcsati kiosztása
%todo elosztó panel dokumentálása

\section{Technikai specifikáció} %todo táblázat
Baud rate: %todo
\par
V+: 12V
%todo áramfelvétel
\par
Vcc: 3.3V
%todo áramfelvétel

\section{Kommunikációs protokoll}
\begin{center}
  \begin{tabular}{ c c }
    Funkció & Küldendő byte \\ \hline \hline
    OldalA panel kiválasztása & F0 \\ \hline
    OldalB panel kiválasztása & F1 \\ \hline
    Verzószám visszaüldése \par A felső hét bit a verzió, az alsó bit az ablak jumperelése & 1101xxxx \\ \hline
    Mindkét ablak fázistolásának párhuzamos írása szubpixelekre vonatkoztatva burstösen & 1110xxxx, ahol xxxx a fázistolás mértéke  \\ \hline
    Szubpixel fényerejének állítása a kiválasztott ablakon & xxxxyyyy, ahol xxxx  a szubpixel sorszáma, yyyy a fényerő \\ \hline
  \end{tabular}
\end{center}

Javaslatok:
\begin{enumerate}
 \item Forráskódolás a kódszókon
 \item 8 bites mélység
\end{enumerate}

\section{Szubpixelek számozása}
%todo

\chapter{MUEB firmware}

\section{Tábla kommunikációs protokoll}
\section{Ethernet protokoll}
Két porton hallgat a MUEB, az egyiken parancsokat lehet neki adni, a másikon pedig animációkat fogad.
\par
Utasítás port: 3000
\par
Animáció port: 2000

\subsection{Parancsok}
\begin{enumerate}
  \item 12V lekapcsolása <-- mátrix végén a csatlakozók nem lesznek feszültség alatt a leszereléskor
  \item 12V bekapcsolása (a véletlen kikapcsolások esetére)
  \item Reboot
  \item Status lekérdezés
  \item Beépített animáció indítása
  \item Beépített animáció leállítása, ekkor az animációk portjára érkezett csomagok jelennek meg.
\end{enumerate}
\subsection{Parancs csomagok felépítése}
A parancs első három bájtja az 'S', 'E' illetve 'M' karakterek ebben a sorrendben, az ezt követő első bájt az utasítás, majd ezt követik az utasítások esetleges argumentumai. 
\subsection{Animáció csomagok felépítése}
Az első bájt szabja meg, hogy melyik ablaknak, és pixelnek szól a csomag. A bájt MSB-je mondja meg, hogy melyik ablak. Ha a bit be van állítva, akkor az kívülről nézve jobb oldali ablak, ha nincs beállítva, akkor kívülről nézve bal ablak. A bájt alsó bitjei határozzák meg a pixelt. Az alsó két bit határozza meg a pixelt. Kívülről nézve bal felső a 0, jobb felső az 1, bal alsó a 2, és jobb alsó a 3. Az első bájt nem használt bitjei 0-ra legyenek állítva. 
\par
A második bájt ábrázolja a pixel vörös komponensét, a harmadik bájt a zöldet, a negyedik bájt a kéket. A skálán a 255 a legfényesebb, a 0 a teljesen kikapcsolt.
\par
A csomagban nem szereplő pixelek változatlan állapotban maradnak.
\subsection{MAC cím}
\section{LED-ek}

\section{Ablak automata}
%\begin{tikzpicture}[shorten >=1pt,node distance=2cm,on grid,auto] 
\begin{tikzpicture}[>=stealth',shorten >=2pt,auto,shorten <=2pt,node distance=5 cm,
                    every state/.style={align=center,minimum size=3cm,text width=2cm},
		    every edge/.style={draw,thick},
		    loop label/.style={draw,align=center,text width=2cm,outer sep=4pt,minimum height=1cm}
		   ]

   \node[state,initial] (vcc 3v3 off)   {3.3V táp kikapcsolva}; 
   \node[state] (vcc 3v3  on) [right=of vcc 3v3 off] {3.3V táp bekapcsolva, de még nincs sikeres kommunikáció}; 
   \node[state] (vcc 12v off) [below=of vcc 3v3 on] {12V-os táp kikapcsolva}; 
   \node[state,accepting](vcc 12v on) [left=of vcc 12v off] {12V-os táp bekapcsolva};
   \path[->] 
          (vcc 3v3 off) edge  node {} (vcc 3v3 on)
          (vcc 3v3 on) edge  node  {üdvözlő üzenet fogadva} (vcc 12v on)
	  (vcc 3v3 on) edge  node  {Timeout} (vcc 3v3 off)
          (vcc 12v on) edge  node  {Parancs} (vcc 12v off)
          (vcc 12v off) edge  node  {Parancs} (vcc 12v on);
\end{tikzpicture}
\par
A MUEB 5 másodperc után bekapcsolja a PWM panel tápellátását, majd vár egy másodpercet az üdvözlőüzenet fogadására. Amennyiben ez nem történik meg, kikapcsolja a csatlakozón a 3.3V-os tápfeszültséget, és 5 másodperc után újra próbálkozik.
\par
Sikeres kommunikációt követően a MUEB bekapcsolja a 12V-ot a táblára. A 12V-ot hálózaton keresztül kapcsolhatjuk ki a megfelelő paranccsal, illetve ebből az állapotból vissza.
\par
Az állapotgép újraindítható szintén egy hálózat felől érkező paranccsal, ekkor a tápellátás megszűnése miatt a kijelző panel is újraindul.

\chapter{Szerver oldali szoftverek}
\chapter{Eszközök karbantartása}
\chapter{Hibakeresés}
\chapter{Egy mátrix felszerelésének menete}
\listoffigures
\chapter{Függelék}
%vicces vagy tanulságos sztorik , szómagyarázat stb.

%todo szószedet

%projektben résztvevő emberek gondolatai, tapasztalatai. a jövő generációinak motiválása

% végső fordításnál kapcsold ki a draft módot !!

\end{document}
